\chapter{Épreuve chronométrée 4}\label{ch:iv1:timed-test-4}

\begin{ivtrains}
La faire d'une seule traite, sans notes ni calculatrice.
\end{ivtrains}

\ivtestsection{A}{6}{40}

\begin{question}\label{iq:iv1:timed-test-4:1}
Un tiroir contient 5 paires de chaussettes différentes. Vous en prenez 4 au hasard. Quelle est la probabilité qu'aucune paire ne soit reconstituée~? Quelle est la probabilité d'exactement une paire reconstituée~?
\end{question}

\begin{question}\label{iq:iv1:timed-test-4:2}
$X$, $Y$ et $Z$ sont des variables exponentielles indépendantes de moyenne 1. Quelle est la probabilité que $X > Y + Z$~?
\end{question}

\begin{question}\label{iq:iv1:timed-test-4:3}
Ajustez une droite par la méthode des moindres carrés aux trois points $(0, 0)$, $(1, 1)$ et $(2, 1)$. Donnez l'ordonnée à l'origine, la pente et le $R^2$.
\end{question}

\begin{question}\label{iq:iv1:timed-test-4:4}
Un jeton part de la case 0, parmi les cases numérotées de 0 à 6. À chaque tour, vous lancez un dé et avancez d'autant de cases, mais si cela vous amène au-delà de la case 6, vous restez où vous êtes. En moyenne, combien de lancers faut-il pour atterrir exactement sur la case 6~?
\end{question}

\begin{question}\label{iq:iv1:timed-test-4:5}
Deux points sont tirés indépendamment et uniformément sur un cercle de rayon 1. Quelle est la distance moyenne en ligne droite entre eux~?
\end{question}

\begin{question}\label{iq:iv1:timed-test-4:6}
Trois joueurs jouent à pierre, papier, ciseaux, en choisissant chacun au hasard. Si les trois joueurs présentent le même geste ou si les trois gestes sont tous différents, ils rejouent. Sinon, les perdants sont éliminés, et s'il en reste deux, ils continuent à jouer. En moyenne, combien de manches faut-il pour désigner un seul vainqueur~?
\end{question}

\ivtestsection{B}{6}{60}

\begin{question}\label{iq:iv1:timed-test-4:7}
Une pièce équilibrée est lancée jusqu'à l'apparition soit de deux piles consécutives, soit de trois faces consécutives. Quelle est la probabilité que les deux piles apparaissent en premier~?
\end{question}

\begin{question}\label{iq:iv1:timed-test-4:8}
Vous tirez des nombres indépendants uniformes sur $[0, 1]$ et vous les multipliez jusqu'à ce que le produit passe sous $\frac12$. Quel est le nombre moyen de tirages~?
\end{question}

\begin{question}\label{iq:iv1:timed-test-4:9}
Un pari est gagnant avec la probabilité $\frac12$ et rapporte 2 pour 1, indépendamment à chaque fois. Quelle fraction de votre richesse faut-il miser à chaque fois pour maximiser la croissance à long terme~? Quel est le taux de croissance par pari~?
\end{question}

\begin{question}\label{iq:iv1:timed-test-4:10}
Deux équipes disputent une série au meilleur de sept matchs, les matchs étant indépendants. Si une équipe gagne chaque match avec la probabilité $0.6$, quelle est la chance qu'elle remporte la série~? Si les équipes sont de force égale, quel est le nombre moyen de matchs~?
\end{question}

\begin{question}\label{iq:iv1:timed-test-4:11}
Deux dés équilibrés sont lancés. Soit $M$ la plus grande valeur et $m$ la plus petite. Quelle est la covariance de $M$ et de $m$~?
\end{question}

\begin{question}\label{iq:iv1:timed-test-4:12}
Le nombre de sinistres en une année suit une loi de Poisson de moyenne 10, et la taille de chaque sinistre est indépendante et suit une loi exponentielle de moyenne 5. Déterminez la moyenne du montant total ainsi que son écart type. Quelle est la chance que le montant total soit nul~?
\end{question}
